\title{Karaage: Japansk fried chicken}
\PrepTime{25}
\CookingTime{}
\CookingTempe{}
\TypeCooking{Kylling}
\NbPerson{4}
\Image{0 100 800 900}{images/karaage.jpg}

\begin{ingredient}
    \begin{main}
        \item Chilisauce\index[ingredient]{Chilisauce}
        \item 2 spsk. Riseddike (el. almindelig eddike)\index[ingredient]{Eddike}
        \item 2 tsk. Hvidløg, revet\index[ingredient]{Hvidløg}
        \item 2 tsk. Ingefær, revet\index[ingredient]{Ingefær}
        \item 400g Kyllingebryst\index[ingredient]{Kyllingebryst}
        \item 4 spsk. Majsstivelse\index[ingredient]{Majsstivelse}
        \item 8 spsk. Hvede-/Kartoffel-mel\index[ingredient]{Mel}
        \item 4dl Friture-/Solsikke-olie\index[ingredient]{Olie}
        \item Ris\index[ingredient]{Ris}
        \item 4 spsk. Soyasauce\index[ingredient]{Soyasauce}
        \item Sukker\index[ingredient]{Sukker}
    \end{main}
\end{ingredient}
\begin{recipe}
    \step{Skær kyllingen i små mundrette bidder og hæld i en skål.}
    \step{Tilføj soyasauce, eddike, hvidløg, sukker og ingefær og lad det marinere i 10--15 minutter.}
    \step{Varm olien op til ca. 160--170 grader i en lille gryde.}
    \step{Bland mel i kyllingemarinaden.}
    \step{Vend hvert stykke kylling i majsstivelse, så det dækker.}
    \step{Friter kyllingen 5--6 stykker af gangen, til de er gyldne og sprøde. Lad dem afdryppe på køkkenrulle.}
    \step{Server med ris og chilisauce.}
\end{recipe}
\begin{notes}

\end{notes}

\setcounter{step}{0}
\setcounter{substep}{0}

\clearpage